% ASBD End-Semester Hackathon report | Jeswin Tobias J | CS23I1013
% Build: cd report && tectonic report.tex
\documentclass[11pt,a4paper]{article}
\usepackage[sfdefault]{plex-sans}         % IBM Plex Sans
\usepackage{plex-mono}                    % IBM Plex Mono for code
\usepackage[a4paper,margin=2.1cm,headheight=14pt]{geometry}
\usepackage{graphicx,booktabs,tabularx,float,xcolor,enumitem,amsmath,array,multirow}
\usepackage[font=small,labelfont=bf]{caption}
\usepackage{fancyhdr,titlesec,tikz}
\usepackage[section]{placeins}   % keep figures inside their own section
\usepackage[colorlinks=true,linkcolor=blue!50!black,urlcolor=blue!50!black]{hyperref}
\graphicspath{{../plots/}}
\newcolumntype{Y}{>{\raggedright\arraybackslash}X}
\newcolumntype{R}{>{\raggedleft\arraybackslash}X}

\definecolor{accent}{HTML}{2a78d6}
\definecolor{boxbg}{HTML}{F3F6FB}
\titleformat{\section}{\Large\bfseries\color{accent}}{\thesection}{0.6em}{}
\titleformat{\subsection}{\large\bfseries}{\thesubsection}{0.6em}{}
\titlespacing*{\section}{0pt}{1.2em}{0.5em}
\titlespacing*{\subsection}{0pt}{0.9em}{0.35em}
\setlist{itemsep=2pt,topsep=3pt,leftmargin=1.3em}
\setlength{\parindent}{0pt}\setlength{\parskip}{5pt}
\emergencystretch=2em

\pagestyle{fancy}\fancyhf{}
\lhead{\small Jeswin Tobias J \,|\, CS23I1013}
\rhead{\small ASBD Hackathon: Tetouan Power Consumption}
\cfoot{\small\thepage}
\renewcommand{\headrulewidth}{0.4pt}

% symbols in maths mode
\newcommand{\ra}{$\rightarrow$\ }
\newcommand{\x}{$\times$}
\newcommand{\gq}{$\geq$}
\newcommand{\AND}{$\wedge$}

% figure followed by its key inference:  \plot{file}{width}{caption}{inference text}
\newlength{\figh}\setlength{\figh}{0.44\textheight}   % cap so two figures fit a page
\newcommand{\plot}[4]{%
\begin{figure}[H]\centering
\includegraphics[width=#2\textwidth,height=\figh,keepaspectratio]{#1}
\caption{#3}\label{fig:#1}\vspace{0.2em}
\colorbox{boxbg}{\parbox{\dimexpr\textwidth-2\fboxsep}{\small\raggedright\textbf{Inference.} #4}}
\end{figure}}
\renewcommand{\topfraction}{0.9}\renewcommand{\bottomfraction}{0.8}
\renewcommand{\textfraction}{0.07}\renewcommand{\floatpagefraction}{0.75}
\setlength{\textfloatsep}{14pt plus 4pt minus 4pt}

\begin{document}

% ------------------------------------------------------------------ title
\hypersetup{pageanchor=false}
\begin{titlepage}
\centering
\vspace*{\fill}
{\Huge\bfseries\color{accent} Power Consumption of Tetouan City\par}
\vspace{0.4cm}
{\Large Descriptive and Predictive Analytics with DIC, DIANA, COFI and PySpark\par}
\vspace{1.4cm}
{\large Name: \textbf{Jeswin Tobias J}\par}
\vspace{2pt}
{\large Roll No: \textbf{CS23I1013}\par}
\vspace{0.8cm}
{\large Analytics \& Systems of Big Data Practice, End-Semester Hackathon\par}
{\large IIITDM Kancheepuram \quad|\quad 21 November 2026\par}
\vspace{1.4cm}
\colorbox{boxbg}{\parbox{\dimexpr\textwidth-2\fboxsep}{\vspace{6pt}\centering
{\small\bfseries\color{accent} PROJECT AT A GLANCE}\par\vspace{6pt}
\small
\begin{tabularx}{0.92\linewidth}{@{}>{\bfseries}p{2.1cm}Y@{}}
Dataset & UCI \#849, Power Consumption of Tetouan City (CC BY 4.0)\\[3pt]
Size & 52,416 readings, one every 10 minutes, 1 Jan to 30 Dec 2017 (364 days)\\[3pt]
Attributes & temperature, humidity, wind speed, two irradiance flows, power of three zones\\[3pt]
Techniques & DIC association rules, DIANA clustering, COFI closed itemsets and rules\\[3pt]
Big data & DIC and COFI on PySpark 4.3 (local mode, 10 cores)\\[3pt]
Code & \url{https://github.com/jeswintobias/asbd-tetouan-power-analytics}\\
\end{tabularx}\vspace{4pt}}}
\vspace*{\fill}
\end{titlepage}
\hypersetup{pageanchor=true}

\tableofcontents
\vspace{1em}
\section*{Report structure}
\begin{tabularx}{\textwidth}{@{}lYl@{}}
\toprule
Section & Content & Plots\\
\midrule
1 & Dataset, assumptions and parameters & none\\
2 & Preprocessing: cleaning, feature engineering, selection, transformation & 01--06\\
3 & Descriptive analytics with 14 plot types & 07--20\\
4 & Association rules with DIC & 21--25\\
5 & Divisive hierarchical clustering with DIANA & 26--30\\
6 & Closed itemsets and closed rules with COFI & 31--34\\
7 & DIC and COFI on PySpark & 35--37\\
8 & Conclusions and limitations & none\\
\bottomrule
\end{tabularx}

% ------------------------------------------------------------------ 1
\section{Dataset, assumptions and parameters}
The dataset gives the electricity demand of three distribution zones in Tetouan, a city in northern Morocco,
along with the local weather, every 10 minutes for 2017. It has no class label. The work is therefore about
describing demand and finding patterns in it: which conditions go with high or low load, and which kinds of day
the year contains.

The question lets us assume parameters as long as we state them. Table~\ref{tab:params} lists every assumption
used anywhere in the report.

\begin{table}[H]\centering\small
\caption{Assumptions and parameters.}\label{tab:params}
\begin{tabularx}{\textwidth}{@{}lYY@{}}
\toprule
Item & Value & Reason\\
\midrule
Validity bounds & temperature $-10$ to 50\,\textdegree C, humidity 0 to 100\,\%, wind \gq 0, irradiance \gq 0, power $>0$ & physical limits\\
Spike filter & median of the 7 readings around each point ($\pm$30 min); thresholds of $\pm$15\,\% for power, $\pm$2\,\textdegree C for temperature, $\pm$10\,\%RH for humidity & the usual 3\x MAD rule flags 13\,\% of normal readings\\
Irradiance floor & values below 1 W/m$^2$ set to 0 & night sensor noise (night maximum 0.154 and 0.271 W/m$^2$)\\
Public holidays & Moroccan national and religious holidays of 2017 (15 days) & day labels\\
Ramadan & 27 May to 25 June 2017 (30 days) & day label\\
Eid periods & 26 June to 2 July, 30 August to 6 September & reading the dips in demand\\
Discretisation & equal-frequency tertiles (Low/Med/High); solar as Dark/Low/High; wind as Calm/Windy split at 3 m/s & each item starts near 33\,\% support\\
Minimum support & 10\,\% (5,242 transactions) & keeps patterns of up to 6 items (768 itemsets)\\
Minimum confidence & 70\,\% & a rule must hold in at least 7 of 10 matching readings\\
DIC station interval & M = 5,000, with 500 to 52,416 tested & about 10\,\% of the data\\
DIANA input & one row per day; 72 values (24 hourly means for each of 3 zones), z-scored per zone, Euclidean distance & distance matrix of 1 MB instead of 22 GB\\
DIANA clusters & k = 5 (k = 2 gives the broad split) & see Section 5.3\\
Spark & PySpark 4.3, \texttt{local[*]}, 10 cores on one laptop; data replicated 1 to 40 times; best of 2 runs & see Section 7\\
\bottomrule
\end{tabularx}
\end{table}

The raw file needed no repair. It has no null cells, no duplicate rows or timestamps and no missing 10-minute
slots, and each of the 364 days has exactly 144 readings. Cleaning was therefore about noise.

% ------------------------------------------------------------------ 2
\section{Preprocessing}
\subsection{Cleaning noisy data}
No reading broke the validity bounds.

We did not use IQR to remove outliers. IQR fences assume a single, stable distribution, and power does not have one.
Its level shifts with the season. All 1,191 readings that IQR flags in Zone 3 fall in July (760) and August
(431), so removing them would delete the summer peak, which is real demand (Figure~\ref{fig:01_iqr_fence_zone3_by_month}).

Spikes were removed with a fixed-threshold rolling-median filter. Each reading is compared with the median
of the 7 readings centred on it. If the gap is larger than the threshold in Table~\ref{tab:params}, the reading is
replaced by that median. We did not use the classic Hampel rule (3\x the scaled MAD). Inside a one-hour window
these series are smooth and the MAD is close to zero, so that rule flags 6,300 to 7,450 readings per column, about
13\,\% of the data. In total the filter replaced 15 power values (4 in Zone 1, 1 in Zone 2, 10 in Zone 3) and
7 weather values (2 temperature, 5 humidity). Irradiance was left alone, since fast changes there come from clouds.

Both irradiance sensors report small positive values at night,
at most 0.154 and 0.271 W/m$^2$. Values under 1 W/m$^2$ were set to 0. This touches 25,182 timestamps (48\,\%),
which looks like a lot, but every changed value was below 1 W/m$^2$ (the largest was 0.998). The changes cover all
readings between 21:00 and 05:59 and none between 09:00 and 17:59, so daylight readings, around 600 W/m$^2$ at noon,
are unaffected.

\plot{01_iqr_fence_zone3_by_month}{0.85}{IQR fences on Zone 3 power, by month.}
{Every IQR outlier in Zone 3 falls in July or August. These are the hottest, busiest weeks of the year, so the
outliers are real summer demand and IQR cannot be used to clean this data.}

\plot{02_hampel_spike_removal}{0.95}{The rolling-median filter on one day of Zone 1.}
{At 12:00 Zone 1 drops from about 37k to 27k kW and stays there. The filter keeps that drop and replaces only the two
isolated readings that fall to about 17k kW. Only 22 values change in the whole year.}

\subsection{Feature engineering}
Every new attribute was added because a later step uses it. ARM in the table means the DIC and COFI transactions.

\begin{table}[H]\centering\small
\caption{Main engineered features.}
\begin{tabularx}{\textwidth}{@{}llYl@{}}
\toprule
Group & Feature & Definition and reason & Used by\\
\midrule
Time & part of day, season & Night 0--6 h, Morning 6--12 h, Afternoon 12--18 h, Evening 18--24 h; seasons by month & ARM, DIANA\\
Time & hour and day of year as sin/cos & keeps 23:00 close to 00:00 & feature check\\
Calendar & weekend, holiday, Ramadan, Eid, day type & non-working and religious days & ARM, DIANA\\
Weather & wind regime & wind sits near 0.08 or near 4.9 m/s, so it is split into Calm and Windy at 3 m/s & ARM\\
Weather & dew point, apparent temp., degree-hours & Magnus and BoM formulas; heating and cooling degree-hours with bases 18 and 24\,\textdegree C & feature check\\
Weather & diffuse fraction, solar mean & cloudiness, and one solar level & ARM, check\\
Weather & hours since sunset & sunset found from irradiance (17:30 in winter, 20:50 in summer) & feature check\\
Power & total, zone shares, load factor & city load, zone mix, daily mean over daily peak & ARM, DIANA\\
Target & load level & total power split into tertiles & feature check\\
\bottomrule
\end{tabularx}
\end{table}

Two candidate features were dropped after checking them. A ``dominant zone'' feature had almost no variance,
because Zone 1 is the largest zone in 99.8\,\% of readings. A zone-transfer feature correlated with load at only
$r = -0.05$. One calendar effect was large enough to keep: between 02:00 and 05:00, load averaged 74.3k kW during
Ramadan and 56.8k kW in the rest of May and June.

\subsection{Selecting task-relevant attributes}
Features with $|r| > 0.9$ were treated as redundant and only one was kept. Relevance was scored by mutual
information with load level. To check the selected set as a whole we used KNN (k = 15, five folds with whole weeks
held out, macro-F1). KNN served only as this check. It is not one of our predictive techniques, since DIANA was
chosen for clustering (Section 5). The check kept the sunset features, which raised macro-F1 from 0.846 to 0.858.
Weekend, Ramadan and Eid flags changed the score by about 0.001, so they stay out of reading-level models and are
used instead as ARM items and as labels for DIANA's days.

\plot{03_feature_correlation_heatmap}{0.8}{Correlation between candidate features.}
{Solar mean tracks general diffuse irradiance, and apparent temperature tracks temperature ($|r|>0.9$), so one of
each pair is dropped. The two irradiance flows are only moderately related and the zones correlate at 0.57 to 0.83,
so both flows and all three zones are kept.}

\plot{04_mutual_information_ranking}{0.8}{Mutual information of each feature with load level.}
{Time of day explains load level best (hour sin 0.40, sunset cos 0.36), followed by solar radiation, temperature and
season. Calendar flags score close to zero because they cover few readings, which is why they are used at day level.}

\subsection{Transformation}
Irradiance is 48\,\% zeros with a long tail, so it was log-transformed (\texttt{log1p}). Features were scaled with
min-max and z-score, and DIANA uses a z-score computed separately for each zone. For DIC and COFI, each 10-minute
reading became one transaction of about 11 items, for example \{Temp=High, Hum=Low, Wind=Calm, Solar=High,
Afternoon, Summer, Workday, Z1=High, \ldots\}. That gives 52,416 transactions over 35 items. For DIANA the readings
were averaged into 364 daily profiles of 72 values each.

\plot{05_log_transform_irradiance}{0.9}{Irradiance before and after the log1p transform.}
{Raw irradiance has a tall spike at zero and a long right tail (skew 1.31). After log1p the tail is compressed,
so a few very bright readings no longer dominate distances.}

\plot{06_item_support}{0.85}{Support of the 35 transaction items.}
{Most items sit between 25 and 35\,\% support, as tertiles should. Holiday (4.1\,\%) and Ramadan (8.2\,\%) are the
rare ones. A 10\,\% minimum support keeps the common patterns but leaves these two items out of the rules, which
we note as a limitation.}

% ------------------------------------------------------------------ 3
\section{Descriptive analytics}
The question asks for at least ten distinct plots from the course. We used 14 types, each picked for the kind of
attribute it suits. Energy totals cover the 364 days in the data.

\begin{table}[H]\centering\small
\caption{Plot types and the question each one answers.}
\begin{tabularx}{\textwidth}{@{}rllY@{}}
\toprule
\# & Plot & Attribute type & Question\\
\midrule
07 & Histogram with KDE & continuous & shape, skew, number of peaks\\
08 & Box plot & numeric by category & demand by day type\\
09 & Violin & numeric by category & Ramadan effect by part of day\\
10 & Line & time series & trend through the year\\
11 & Stacked area & time series, part of a whole & each zone's share over the day\\
12 & Scatter & numeric vs numeric & temperature and load\\
13 & Donut & part of a whole & energy share per zone\\
14 & Grouped bar & category by category & season and zone\\
15 & Pareto & ranked contribution & which hours carry the energy\\
16 & Radar & profile on several metrics & how the zones differ\\
17 & Heatmap & two time axes & when the peak happens\\
18 & Treemap & hierarchy & energy by season, then zone\\
19 & Dot plot & small summary (36 values) & monthly level per zone\\
20 & Stem-and-leaf & small summary (52 values) & weekly load with every value kept\\
\bottomrule
\end{tabularx}
\end{table}

\plot{07_histogram_kde}{0.95}{Histogram with KDE for every numeric attribute.}
{Temperature is close to symmetric (skew 0.20) and humidity leans left ($-$0.63). Wind has two peaks: 61\,\% of
readings are under 0.5 m/s and 39\,\% are over 4 m/s, with little in between, so it is treated as two regimes.
Irradiance is 48\,\% zeros, which led to the log transform, and Zone 3's long right tail is its summer rise.}

\plot{08_box_zone_by_daytype}{0.85}{Box plot of zone demand by day type.}
{At weekends the median falls 9.5\,\% in Zone 2 and 4.5\,\% in Zone 1, while Zone 3 rises 1.2\,\%. Zone 2 depends
most on working days. Zone 3 is 7.3\,\% higher on holidays, but 7 of the 15 holidays fall in July to September, its
busy season, so this is probably a seasonal effect.}

\plot{09_violin_ramadan_part_of_day}{0.85}{Violin plot of city load by part of day, Ramadan against the rest of May and June.}
{The comparison uses the rest of May and June, which had similar weather. During Ramadan the median night load is
41\,\% higher (79.5k against 56.3k kW), the morning is 11\,\% lower and the evening 8\,\% higher. Demand moves from
morning to night. Pre-dawn meals and later working hours could explain this, though the data does not record either.}

\plot{10_line_daily_trend}{0.95}{Daily mean consumption through 2017.}
{All three zones peak in July and August. Zone 3 swings the most: its highest month is 2.55 times its lowest, against
1.26 for Zone 1 and 1.40 for Zone 2. After summer the zones part ways, with Zone 3 down to 11.0k kW in December and
Zone 2 up to 23.7k kW. The sharp dips at the end of June and around 1 September fall on Eid al-Fitr and Eid al-Adha.}

\plot{11_stacked_area_daily_profile}{0.85}{Stacked area chart of the average day, split by zone.}
{City load is lowest at 06:50 (49.7k kW) and highest at 20:10 (98.4k kW), about twice as much. Zone 3's share grows
from 23.5\,\% at the low point to 26.5\,\% at the peak, and its own daily swing (2.26 times) is the largest. The
evening climb lines up with sunset, so lighting is one likely cause.}

\plot{12_scatter_temp_vs_load}{0.8}{Temperature against city load, coloured by hour.}
{Across all readings $r = 0.49$, but within a single hour the average is $r = 0.59$ and reaches 0.77, because time
of day hides part of the temperature effect. Readings above 30\,\textdegree C average 95.5k kW against 71.2k kW
overall, 34\,\% more. Evening readings stay high even at 10 to 15\,\textdegree C.}

\plot{13_donut_energy_share}{0.6}{Each zone's share of energy over the 364 days.}
{The city used about 622 GWh over the 364 days. Zone 1 took 45.4\,\% (283 GWh), Zone 2 29.5\,\% (184 GWh) and
Zone 3 25.0\,\% (156 GWh). Zone 1 is close to half the city, while Zone 3 is the smallest and the least steady
(see plots 10 and 16).}

\plot{14_grouped_bar_season_zone}{0.8}{Mean load by season and zone.}
{From winter to summer, Zone 3 rises 59.3\,\%, Zone 1 17.5\,\% and Zone 2 12.3\,\%. Zone 3 is lowest in autumn
(13.7k kW) and passes Zone 2 only in summer, so the order of the zones depends on the season. That is why season is
one of the ARM items.}

\plot{15_pareto_energy_by_hour}{0.85}{Pareto chart of energy by hour of day.}
{The five hours from 18:00 to 22:00 carry about 27\,\% of each day's energy, and 20:00 alone carries 5.74\,\%, twice
as much as 06:00. Reaching 80\,\% of the energy takes 18 of the 24 hours, so there is no 80/20 pattern; the base load
is large. Any effort to cut the peak should focus on the evening.}

\plot{16_radar_zone_profiles}{0.6}{Radar chart of the zones on six metrics, each scaled by its maximum.}
{Zone 1 is the largest and steadiest (coefficient of variation 0.22). Zone 3 varies most (0.37) and is 1.59 times
higher in summer than in winter. Zone 2 has the lowest weekend-to-workday ratio (0.92). All three have an
evening-to-night ratio near 1.8, so the evening peak is city-wide. Zone 3's seasonal pattern fits a coastal area
with summer visitors, but the data cannot confirm that.}

\plot{17_heatmap_hour_month}{0.72}{Mean city load by hour and month.}
{The peak hour moves with sunset, from 18:00 in November and December to 20:00 from April to August. The highest
cell is July at 20:00 (121k kW) and the lowest is May at 05:00 (45.6k kW). Between midnight and 03:00, June reads
74k to 88k kW against 58k to 68k kW in May, which is Ramadan. July and August stay high through the whole day.}

\plot{18_treemap_season_zone}{0.72}{Treemap of energy by season, then zone (area in GWh).}
{Summer is the largest season with 184 GWh (29.6\,\%) and winter the smallest with 142 GWh. Summer in Zone 1 is the
largest single block (12.6\,\% of the total) and autumn in Zone 3 the smallest (4.8\,\%). Summer is the only season
where Zone 3 (54 GWh) uses more than Zone 2 (51 GWh).}

\plot{19_dot_plot_monthly_zone}{0.85}{Monthly mean per zone (36 values).}
{Zone 1 is always on top, between 29k and 36k kW. Zone 3 is above Zone 2 in April and July and equal in August.
From September the gap opens until Zone 2 is 12.6k kW ahead in December. Zones 2 and 3 move in opposite directions
through the year, and DIANA picks up the same swap later.}

The weekly summary has only 52 values, which is why we used a stem-and-leaf plot for it. Each weekly value can be
read off the plot to the nearest 1,000 kW, so the summary could be rebuilt from the figure, much like reading a
short CSV, and the shape of the distribution is still visible. A histogram would hide the values inside its bins.

\plot{20_stem_and_leaf_weekly_load}{0.6}{Stem-and-leaf plot of weekly mean city load (52 weeks). Key: 7L\,|\,3 = 73,000 kW.}
{37 of the 52 weeks fall between 63k and 74k kW (median 68k, quartiles 67k and 72k). All 8 weeks at or above 80k kW
start between 2 July and 20 August, and the highest is 95k. The right tail is the summer peak.}

These plots point to time of day, season, Ramadan and day type as the main structure in the data. Those become the
main items for DIC and COFI, and the labels used to read DIANA's clusters.

% ------------------------------------------------------------------ 4
\section{Association rules with DIC}
\subsection{Method}
Apriori reads the whole dataset once for each itemset size. DIC (Dynamic Itemset Counting) puts a station every
M transactions and lets a larger itemset start counting at any station, as soon as all of its subsets look frequent.
Each itemset is in one of four states. A dashed circle is still counting and below minimum support. A dashed square
is still counting but already above it. A solid square has finished and is frequent, and a solid circle has
finished and is not. An itemset stops counting once it has seen every transaction, one full lap after it started.
We wrote DIC from scratch in \texttt{src/arm.py}.


\subsection{Results}
With 10\,\% support, 70\,\% confidence and M = 5,000, DIC found 768 frequent itemsets (33, 224, 281, 172, 52 and 6
of sizes 1 to 6). They match mlxtend's Apriori and our own Apriori in both itemsets and counts. DIC needed 5.57
passes over the data at M = 5,000 and 4.92 at M = 500, against 6.0 for Apriori, so up to 18\,\% fewer scans. In
memory, though, DIC was slower (0.63 s against 0.18 s) and held more counters at once (831 against 686). It helps
when each scan is expensive, as with data on disk or spread across machines.

1,800 rules reached 70\,\% confidence. We removed a rule when a smaller left-hand side gave the same right-hand side
with equal or higher confidence, which left 1,098 rules, all with lift above 1. Of these, 43 have only conditions
(time, weather, day) on the left and only load levels on the right. Those are the rules that say when demand is high
or low.

\begin{table}[H]\centering\small
\caption{Strongest rules from conditions to load (all rules: \texttt{outputs/dic\_rules.csv}).}
\begin{tabular}{@{}lrrr@{}}
\toprule
Rule & Support & Confidence & Lift\\
\midrule
Night \AND Temp=Low \ra Z1=Low & 10.5\,\% & 96.8\,\% & 2.90\\
Night \AND Temp=Low \ra Total=Low & 10.2\,\% & 94.6\,\% & 2.84\\
Night \AND Wind=Calm \ra Z1=Low & 13.6\,\% & 86.0\,\% & 2.58\\
Summer \AND Temp=High \AND Wind=Windy \ra Z3=High & 14.6\,\% & 86.0\,\% & 2.58\\
Evening \AND Workday \ra Total=High & 14.6\,\% & 84.8\,\% & 2.54\\
Evening \ra Total=High & 21.1\,\% & 84.4\,\% & 2.53\\
\midrule
Temp=High \AND Z1=High \AND Z3=High \ra Summer \AND Total=High \AND Z2=High & 10.3\,\% & 75.7\,\% & 5.96\\
\bottomrule
\end{tabular}
\end{table}
These rules describe readings taken at the same moment. They say which conditions and load levels occur together,
not what load will be later.

\plot{21_dic_itemset_states}{0.9}{Number of itemsets in each DIC state at every station (M = 5,000).}
{Pairs start counting at the second station, after 10,000 transactions, long before the first pass ends. The number
of dashed circles peaks at about 720 around pass 1.4, and everything has finished after 5.57 passes. Holiday
(4.1\,\%) and Ramadan (8.2\,\%) end as solid circles.}

\plot{22_dic_passes_time_vs_M}{0.95}{Effect of the station interval M and of data order on DIC.}
{Smaller M means fewer passes (4.92 at M = 500). With M equal to the whole dataset DIC behaves like Apriori and
needs 6 passes. More stations also mean more bookkeeping, and the run takes 4.2 s at M = 500. Shuffled data needs
fewer passes at medium M (5.19 against 5.57), because readings in time order arrive season by season, and DIC works
best when the data is evenly mixed.}

\plot{23_dic_rules_support_confidence}{0.8}{Support against confidence for the 1,098 rules, coloured by lift.}
{Most rules sit between 10 and 15\,\% support. The highest lifts, around 5 to 6, all describe the same situation:
summer, high temperature and high load in every zone. In hot weather the three zones are high together. Broad,
high-support rules are closer to what chance alone would give.}

\plot{24_dic_top_context_load_rules}{0.9}{Top rules from conditions to load, by confidence.}
{Cold nights go with low load 95 to 97\,\% of the time, evenings with high load 84\,\% of the time, and hot, windy
summer readings with high Zone 3 load 86\,\% of the time.}

\plot{25_dic_lift_matrix_context_load}{0.85}{Lift for each pair of condition and load level.}
{Evening has a lift of 2.53 with high total load and almost none (0.01) with low load. Ramadan has a lift of 1.52
with high load, but its support of 8.2\,\% is under the 10\,\% threshold, so it cannot appear in a rule. Holiday
sits near 1.0 and weekend at 1.11 with low load, so neither changes city demand much.}

% ------------------------------------------------------------------ 5
\section{Divisive hierarchical clustering with DIANA}
\subsection{Why DIANA and not KNN}
The question allows KNN or DIANA, whichever suits the data. This data has no class label and the goal is to find
types of day, so an unsupervised method fits. KNN is a supervised classifier and would need a label made up from
the power columns. DIANA also gives the dendrogram the question asks for.

\subsection{Method}
DIANA starts with every object in one cluster and keeps splitting. At each step it takes the cluster with the
largest diameter. The object with the largest mean distance to the rest leaves first and starts a splinter group.
Other objects follow while they are closer, on average, to the splinter group than to the old cluster. This repeats
until every object is on its own, and the diameter at each split gives the height in the dendrogram. We wrote it
from scratch in \texttt{src/diana.py}.

Each object is one day, 364 in all. With one row per 10-minute reading the distance matrix would need about 22 GB,
with hourly rows about 0.6 GB, and with daily rows about 1 MB. Each day is described by 72 values, the hourly mean of
each zone, because two days can have the same city total while different zones peak. Each zone's values are
z-scored separately, or Zone 1 (around 32k kW) would outweigh Zone 3 (around 18k kW). Season, day type, Ramadan and
Eid are not given to DIANA. We used them only afterwards, to read the clusters.


\subsection{Choosing k and checking quality}
The tree has a divisive coefficient of 0.932 and a cophenetic correlation of 0.817. Two clusters give the clearest
split (silhouette 0.491): 73 hot-season days against the rest of the year. We chose five clusters to describe the
year in more detail. Silhouette has a local peak there (0.387), and the split height drops after the fourth split,
from 11.95 to 8.83. The peak is narrow, since k = 6 scores 0.382 and k = 7 scores 0.376, so k = 5 is a descriptive
choice and k = 6 or 7 would also be defensible. Clustering the raw kW values gives nearly the same partition
(adjusted Rand index 0.854).

\plot{26_diana_choose_k}{0.95}{Silhouette and split height used to choose the number of clusters.}
{Silhouette is 0.491 at k = 2 and has a smaller local peak of 0.387 at k = 5, with 0.382 at k = 6 and 0.376 at k = 7.
Split height falls after the fourth split. Two clusters give the broad hot-season split, and five give day types
that can be described.}

\plot{27_diana_dendrogram}{1.0}{DIANA dendrogram of the 364 days, cut at k = 5 (colours show the clusters).}
{The first split, at a diameter of 19.8, separates the 73 hot-season days from the rest, so season is the main
structure in the data. The two Ramadan-associated groups are more spread out inside than the other three.}

\subsection{The five day types}
\begin{table}[H]\centering\small
\caption{Day types found by DIANA (k = 5). Calendar labels were not used to form them.}
\begin{tabularx}{\textwidth}{@{}llrY@{}}
\toprule
Cluster & Name & Days & Profile\\
\midrule
C0 & Base & 180 & January to mid-October, mild (16.6\,\textdegree C), 67.7k kW, peak at 20:00; includes the Eid al-Adha week\\
C1 & Ramadan-associated A & 24 & 18 Ramadan days (27 May to 12 June), the Eid al-Fitr days and 4 others; Zone 1 at 03:00 is 32.1k kW against 23.7k in Base\\
C2 & Ramadan-associated B & 21 & 12 Ramadan days (13 to 25 June), the week after Eid and a few summer days; same night pattern, warmer (23.3\,\textdegree C)\\
C3 & Peak summer & 52 & 4 July to 25 August, 26.7\,\textdegree C, 89.0k kW (31\,\% above Base); Zone 3 at 20:00 is 40.0k kW\\
C4 & Autumn and winter & 87 & mid-September to December; Zone 3 falls to 17.6k kW at 20:00 while Zone 2 rises to 31.2k; peak at 18:00\\
\bottomrule
\end{tabularx}
\end{table}

All 30 Ramadan days ended up in C1 and C2, even though DIANA never saw the calendar. The two clusters also hold
15 other days, so they are better described as Ramadan-associated than as Ramadan clusters. Pre-dawn meals,
travel around Eid and early-summer visitors in Zone 3 could all play a part; we did not test these. Weekends never
form a cluster of their own: every cluster is 21 to 38\,\% weekend days.

\plot{28_diana_cluster_profiles}{1.0}{Average 24-hour profile of each cluster, per zone.}
{The Ramadan-associated clusters have high night load. Peak summer is high all day, with Zone 3 at 40k kW in the
evening. In autumn and winter Zone 2 passes Zone 3 and the peak moves to 18:00. A profile of city totals would hide
this swap, which is why each zone keeps its own 24 columns.}

\plot{29_diana_calendar}{0.85}{The 2017 calendar coloured by DIANA cluster.}
{The clusters form unbroken stretches of the year. C1 and C2 cover late May to early July around Ramadan and C3
covers July and August. DIANA found these seasons without being given any dates.}

\plot{30_diana_silhouette}{0.85}{Silhouette of every day, grouped by cluster.}
{Peak summer (0.467), autumn and winter (0.449) and Base (0.398) are well separated. The Ramadan-associated clusters
are weak (0.167 and 0.082) and contain all 8 days with a negative score. Those days share the night pattern but their
daytime load sits between Base and summer. The next split, at k = 6, separates exactly the two Eid al-Fitr days
(26 and 27 June).}

% ------------------------------------------------------------------ 6
\section{Closed itemsets and closed rules with COFI}
\subsection{Method}
COFI mines frequent itemsets with one small tree per item. We wrote it from scratch in \texttt{src/cofi.py}. It
works in four steps.
\begin{enumerate}
\item Merge identical transactions (52,416 become 3,986 distinct ones with counts), count item supports and build an
FP-tree with items in order of support.
\item Take the items one at a time, least frequent first. For item $i$, build a Co-Occurrence Frequent-Item tree from
the FP-tree paths above $i$. Each node keeps a count and a participation counter.
\item Mine that tree. A branch's count is the node count minus its participation, every sub-pattern of the branch
(with $i$) gets that count, and participation is then raised along the branch so nothing is counted twice.
\item Throw the tree away before building the next, so only one COFI-tree is in memory at a time.
\end{enumerate}
An itemset is closed if no superset one item larger has the same support, and maximal if no such superset is
frequent. Every maximal itemset is closed and every closed itemset is frequent. Closed rules use the Min-Max bases.
Exact rules go from a minimal generator $G$ to the rest of its closure, with 100\,\% confidence. Approximate rules
go from $G$ to a larger closed set $Z$, and the set is reduced so no rule is implied by another.


\subsection{Results}
At 10\,\% support and 70\,\% confidence COFI finds the same 768 frequent itemsets and counts as DIC and Apriori. The
FP-tree has 14,275 nodes, and the largest COFI-tree (for \emph{Afternoon}) has 429, 3\,\% of the FP-tree.

753 of the 768 itemsets are closed (98\,\%) and 198 are maximal (26\,\%). The full list, with minimal generators,
is in \texttt{outputs/cofi\_closed\_itemsets.csv}, and the supports of all 768 frequent itemsets can be rebuilt
from the closed ones. The 15 itemsets that are not closed all come from arithmetic: if Z1, Z2 and Z3 are all High,
Total is High too (and the same for Low), because Total = Z1 + Z2 + Z3.

The Min-Max basis has 15 exact and 1,111 approximate rules, 1,126 in all, against 1,800 rules from the frequent
itemsets (37\,\% fewer). We checked that the support and confidence of all 1,800 rules can be recomputed from the
753 closed itemsets; the largest error was $2.2\times10^{-16}$. In theory the full rule set can also be derived from
the basis itself, but we did not implement that, so we only claim the rebuild from the closed itemsets.

\begin{table}[H]\centering\small
\caption{Longest closed itemsets (6 items, all maximal).}
\begin{tabular}{@{}lrr@{}}
\toprule
Closed itemset & Count & Support\\
\midrule
\{Night, Solar=Dark, Total=Low, Wind=Calm, Z1=Low, Z2=Low\} & 6,270 & 12.0\,\%\\
\{Night, Solar=Dark, Total=Low, Workday, Z1=Low, Z2=Low\} & 6,027 & 11.5\,\%\\
\{Solar=Dark, Temp=Low, Total=Low, Wind=Calm, Z1=Low, Z2=Low\} & 5,428 & 10.4\,\%\\
\{Summer, Temp=High, Total=High, Z1=High, Z2=High, Z3=High\} & 5,411 & 10.3\,\%\\
\bottomrule
\end{tabular}
\end{table}
An exact closed rule is \{Z1=High, Z2=High, Z3=High\} \ra \{Total=High\} with 100\,\% confidence, which is the sum
again. An approximate one is \{Summer, Z1=High, Z2=High, Z3=High\} \ra \{Temp=High, Total=High\} with 88.6\,\%
confidence and lift 5.1. The strongest closed rules from conditions to load are the same as DIC's: cold nights go
with low load (95\,\%), hot windy summers with high Zone 3 load (86\,\%) and evenings with high load (84\,\%).

\plot{31_cofi_tree_sizes}{0.85}{Size of each item's COFI-tree compared with the FP-tree.}
{The largest tree, for Afternoon, has 429 nodes, 3\,\% of the 14,275-node FP-tree. The most frequent item, Workday,
has an empty tree because no item sits above it. Since only one small tree exists at a time, mining needs very
little memory.}

\plot{32_fi_cfi_mfi_vs_support}{0.9}{Frequent, closed and maximal itemsets for minimum support from 2 to 20\,\%.}
{At 20\,\% support 99\,\% of frequent itemsets are closed. At 2\,\% the share falls to 85\,\% (15,924 of 18,770),
because lower support lets in long itemsets in which exact co-occurrence is more common. Maximal itemsets are always
far fewer (12 to 46\,\%), but they do not keep support counts.}

\plot{33_rule_set_sizes}{0.8}{Number of rules at 70\,\% confidence under each method.}
{The 1,800 rules drop to 1,098 with DIC's pruning and to 1,126 in the Min-Max basis (15 exact, 1,111 approximate).
The basis is about as small as the pruned set, and its closed itemsets reproduce every one of the 1,800 rules.}

\plot{34_cofi_dic_apriori_runtime}{0.85}{Run time of COFI, DIC and Apriori against minimum support (log scale).}
{All three were timed from the raw transactions, best of two runs. DIC is the slowest at every support up to 10\,\%
(0.37 s against 0.06 s for COFI at 10\,\%). COFI beats Apriori between 8 and 20\,\% (0.06 s against 0.14 s at
10\,\%) and loses at 5\,\% (0.74 s against 0.57 s). COFI still enumerates candidates, 118,045 sub-patterns at 10\,\%
and 4.3 million at 5\,\%, against Apriori's 1,575 and 6,453. What it avoids is one global candidate set and repeated
scans of the data, and the growing number of sub-patterns is why it slows down at low support.}

% ------------------------------------------------------------------ 7
\section{DIC and COFI on PySpark}
\subsection{How the two techniques run on Spark}
For DIC, the driver keeps the same state machine as the single-machine code. Each station's block of M transactions
is stored as 10 partitions and counted in parallel with \texttt{sc.runJob}. Transactions are stored as 64-bit masks,
so checking whether a transaction contains an itemset is a single AND.

For COFI, each item's tree is independent of the others, so each becomes one Spark task that runs the same mining code
as Section 6. Before that, a map-side combiner merges identical transactions inside each partition, much like a
partial \texttt{GROUP BY} in SQL, so only about 4,000 weighted transactions are turned into FP-paths. DIANA was not
moved to Spark, because each split depends on the one before and its input is only 1 MB.

The tests ran on one laptop with Spark in \texttt{local[*]} mode on 10 cores. To get larger inputs we repeated the
52,416 transactions 1, 5, 10, 20 and 40 times, up to 2,096,640. Supports do not change under repetition, so every
count must be exactly k times the original. This held at all five sizes for both techniques. Times are the best of
two runs. Mining times use data already cached in Spark, and loading and start-up are listed separately. The
single-machine COFI merges duplicates in the same way, so Spark does not get a shortcut the baseline lacks.

\subsection{Results}
\begin{table}[H]\centering\small
\caption{Performance at 40 times the data (2,096,640 transactions), not counting Spark start-up (2.7 s).}
\begin{tabularx}{\textwidth}{@{}lRRRRR@{}}
\toprule
Algorithm & Python (s) & Spark mining, cached (s) & Speed-up, cached & Spark load and mining (s) & Speed-up with loading\\
\midrule
DIC & 3.659 & 1.219 & 3.00\x & 2.282 & 1.60\x\\
COFI & 0.477 & 0.369 & 1.30\x & 0.707 & 0.68\x\\
\bottomrule
\end{tabularx}
\end{table}

For DIC, Python's time grows in step with the data (0.33 to 3.66 s) while Spark stays near 1.2 s. Spark catches up at
about 0.55 million transactions and is 3.0 times faster at 2.1 million, or 1.6 times once loading is counted.

For COFI, Spark catches up at about 1.5 million transactions and is 1.30 times faster at 2.1 million on cached data.
Once the time to load the data into Spark is added, it is slower than Python (0.68 times).

The speed-up grew across the sizes we tested. Each extra million rows cost Python 1.65 s for DIC and 0.21 s for COFI,
and Spark 0.08 s for each. We only measured this on repeated data, so we cannot say the same holds for data with many
more distinct transactions. On the original 52k transactions Python wins for both (Spark runs at 0.27 times the
speed for DIC and 0.23 for COFI), because start-up and per-job costs dominate.

\plot{35_spark_runtime_vs_size}{1.0}{Run time against data size: Python on 1 core and PySpark on 10 cores.}
{Python's DIC time grows in a straight line while Spark's stays near 1.2 s. For COFI, Python goes from 0.05 to 0.48 s
and Spark from 0.20 to 0.37 s. Most of Spark's cost is fixed, while Python pays for every row, so Spark pulls ahead
as the data grows.}

\plot{36_spark_speedup_vs_size}{0.85}{Speed-up of PySpark over single-machine Python (above 1 means Spark is faster).}
{DIC passes the break-even line near 0.55 million transactions and reaches 3.0 times; COFI passes it near 1.5 million
and reaches 1.30 times. Spark only helps above roughly half a million to 1.5 million transactions, and on the original
data it is slower.}

\subsection{Where the COFI gain comes from}
\begin{table}[H]\centering\small
\caption{Naive and combiner versions of Spark COFI: same data (2.1 million), same 10 cores, same output.}
\begin{tabularx}{\textwidth}{@{}lRRRRR@{}}
\toprule
Version & Spark jobs & Rows after merging & FP-path records & Time (s) & Speed-up over Python\\
\midrule
Naive & 3 & 2,096,640 & 22,976,640 & 2.54 & 0.19\x\\
Combiner & 2 & 79,720 & 43,508 & 0.38 & 1.25\x\\
\bottomrule
\end{tabularx}
\end{table}
The first Spark version turned every row into FP-paths before merging anything. Merging identical transactions
first, inside each partition, cut the records generated by a factor of 528 and made Spark COFI 6.7 times faster. On
Spark most of the time went into creating and moving records; the mining itself was cheap. We counted records generated; we did not
measure bytes sent over the network.

We also tried a version that first filters each transaction to its frequent items and then emits every prefix. It
gave the same output but took 0.94 to 1.44 s at 40 times the data, because it reads the raw rows twice. One point in
its favour is that it never collects the distinct transactions on the driver. Our version does, which is fine with
about 4,000 of them but would become a bottleneck with millions.

\plot{37_spark_speedup_vs_cores}{0.8}{Speed-up against number of cores for 2.1 million transactions.}
{On 10 cores COFI runs 3.8 times faster than on one (1.39 to 0.36 s) and DIC 2.4 times faster (2.93 to 1.21 s).
Neither is close to linear. Solving Amdahl's law, $S = 1/(f + (1-f)/p)$, for the serial part gives about 18\,\% for
COFI (job scheduling and driver work) and about 35\,\% for DIC (the state machine on the driver and one job per
station).}

Spark start-up (2.7 s) and loading (0.3 to 1.1 s at 40 times the data) are paid once. With start-up included, Spark
DIC at 2.1 million transactions takes about 5.0 s against 3.7 s in Python, so a single run at this size does not pay
off. Spark becomes worthwhile when the session and cached data are reused, for example to mine at several support
levels, or when the data no longer fits on one machine.

% ------------------------------------------------------------------ 8
\section{Conclusions and limitations}
\begin{table}[H]\centering\small
\caption{Summary of findings.}
\begin{tabularx}{\textwidth}{@{}lY@{}}
\toprule
Part & Main result\\
\midrule
Preprocessing & complete data; IQR rejected because it would delete real summer demand; 22 spike values fixed; night irradiance set to zero; 35-item transactions; 364 daily profiles of 72 values\\
Descriptive & evening peak about twice the morning low and moving with sunset; Zone 3 most seasonal (59\,\% higher in summer); Ramadan moves load to the night (41\,\% higher); Zone 2 most tied to working days\\
DIC & 768 itemsets in 5.57 passes against Apriori's 6.0; 1,098 rules after pruning; cold nights with low load (95\,\%), evenings with high load (84\,\%), hot windy summers with high Zone 3 load (86\,\%)\\
DIANA & five day types (divisive coefficient 0.932); season is the first split; all 30 Ramadan days fall into two clusters without the calendar; Zones 2 and 3 swap places after summer\\
COFI & 753 closed and 198 maximal itemsets; 1,126 Min-Max rules; every rule can be recomputed from the closed itemsets\\
PySpark & identical results at every size; on cached data 3.0 times faster for DIC and 1.30 times for COFI at 2.1 million transactions; the combiner made Spark COFI 6.7 times faster\\
\bottomrule
\end{tabularx}
\end{table}

The work has some limits. Holiday (4.1\,\%) and Ramadan (8.2\,\%) are below the 10\,\% support threshold, so no
rule mentions them; a lower threshold or mining those days separately would be needed. The two Ramadan-associated
clusters are weak (silhouette 0.17 and 0.08), include 15 days outside Ramadan, and k = 5 is only slightly ahead of
k = 6 and 7. Causes such as tourism, lighting or meal times are not in the data, so we give them only as possible
explanations. Finally, Spark was tested on one machine with repeated data, so the number of distinct transactions
never grew; a real cluster or more varied data could behave differently.

\end{document}
